\documentclass[11pt]{article}
\usepackage{amsmath,amssymb}
\usepackage[margin=1in]{geometry}
\usepackage{hyperref}

\title{Ideas Parked --- Not Yet Merged (v15)}
\author{Staging document; each entry marks its eventual destination}
\date{}

\begin{document}
\maketitle

\noindent Two conversations produced one finding, not two. A
density/energy pair built from the same per-vector quantity
($S_{\text{atom}}$ and $\mathcal{E}_{\text{atom}}$, both sums of
$r_i$ transformed) turned out to be forced into a one-directional
bound. A second pair, density redefined as bare count and energy
redefined as a shell-weighted strength distribution, turned out to
have no relationship at all, in either direction. This entry states
the rule both results are instances of, gives a reusable test for
any future pair of atom-level metrics, and sketches where the
distinction might actually matter --- all of it exploratory, none
of it adopted.

\section{The rule}

Call a metric on a Collapsed atom \emph{separable} if it is a sum
over the atom's own vectors of some fixed per-vector function:
$M=\sum_{i=1}^N f(r_i)$. Every metric discussed so far is of this
form, including the ones that do not look it: $N$ itself is
$\sum_i 1$, the constant function.

\paragraph{The test.} For two separable metrics $M_1=\sum f(r_i)$
and $M_2=\sum g(r_i)$: if $f(r)\le g(r)$ for every $r\in[0,1]$, then
$M_1\le M_2$ for every atom, with no exceptions --- summing a
pointwise inequality preserves it. If neither $f\le g$ nor $f\ge g$
holds everywhere on $[0,1]$ --- if the two functions cross --- no
such bound exists in either direction, and the two metrics are free
to move independently, as far apart as the atom's own vectors allow.

\paragraph{Checked against both findings on file.} $S_{\text{atom}}$
and $\mathcal{E}_{\text{atom}}$ use $f(r)=r$ and $g(r)=r^2$; checked
directly, $r^2\le r$ holds across all of $[0,1]$, so
$\mathcal{E}_{\text{atom}}\le S_{\text{atom}}$ is forced, matching
the identity already found. $N$ and the shell-weighted energy
sketch use $f(r)=1$ and $g(r)=\lfloor r/w\rfloor\cdot r$; checked
directly, neither $f\le g$ nor $f\ge g$ holds on all of $[0,1]$ (the
constant function sits below $g$ near $r=1$ and above it near
$r=0$), so no bound exists, matching the unbounded result already
found. One test, both results, no coincidence.

\paragraph{Why this is worth having as a standing test.} Any future
proposal for a new atom-level metric can be checked against every
metric already on file in one step --- write both as $f$ and $g$ on
$[0,1]$, check whether the two curves cross --- rather than
re-running numerical trials each time a new pair is proposed. A
crossing means the pair is informative together (two independent
readings of the atom); no crossing means one is redundant given the
other, except for the gap between them, which is its own quantity
(below).

\section{The gap, where one exists, is itself a quantity}

Where a bound holds, $M_2-M_1=\sum\bigl(g(r_i)-f(r_i)\bigr)$ is not
nothing --- it is a separable metric in its own right, built from the
per-vector function $g-f$. For $S_{\text{atom}}-\mathcal{E}_{\text{atom}}$,
that function is $r-r^2$, maximized at $r=0.5$ and zero at both
$r=0$ and $r=1$ (already found). The general pattern: wherever two
metrics are forced into a bound, the slack between them measures
something like efficiency --- how much of one quantity a vector
contributes without contributing a matching share of the other ---
and is itself worth naming once a specific bounded pair is adopted,
not only in this one case.

\section{Where this might be used, if any of it is adopted}

None of this is proposed as settled; each is a direction worth
testing, not a result.

\paragraph{A minimal, non-redundant reporting set.} With seven
metrics already on file and at least two more count-based ones
proposed, not all of them are worth reporting side by side. The test
above sorts the existing inventory into two kinds: pairs that are
bounded (reporting both is partly redundant --- the gap is the real
second number) and pairs that are free (reporting both is two
genuinely independent readings). A table applying the test to every
pair already in the inventory --- $S_{\text{atom}}$,
$\mathcal{V}_{\text{atom}}$, $\gamma_{\text{atom}}$,
$\mathcal{E}_{\text{atom}}$, $N$, $D$ --- would say, once, which
combinations are worth stating together.

\paragraph{A two-axis placement for comparing atoms.} Where a pair is
free (density and shell-energy, as checked here), plotting a
population of Collapsed atoms on those two axes is a meaningful
scatter, not a degenerate one --- nothing forces the points onto or
under a curve the way a bounded pair would. Where a pair is bounded
instead ($S_{\text{atom}}$ against $\mathcal{E}_{\text{atom}}$),
the same plot is still informative, but differently: every atom
must fall on or below the line $\mathcal{E}=S$, and \emph{how far
below} is the efficiency reading from Section~2, not a free
coordinate.

\paragraph{A possible reading for Inference's own evidence
questions, offered cautiously.} The common split in evidence
aggregation --- many weak signals versus one strong one --- is
roughly what a free density/energy pair would distinguish: an atom
that is dense but low-energy under the shell reading looks like many
weak contributions; one that is light but high-energy looks like a
single strong one. This is an analogy to be tested against
Inference's own existing reference-atom mechanism, not a replacement
for it --- it may turn out to duplicate a distinction that mechanism
already makes some other way, which would need checking before
anything is built on it.

\paragraph{A post-hoc description of Assembly's frozen camps, not a
new input to them.} Section~5 of ideas\_v13 closed the question of
feeding any density- or mass-like weighting into Assembly's own
motion step, on direct, tested evidence that it breaks monotonic
convergence. That finding is not reopened here. What is not closed:
describing a camp's own resulting density and energy, once frozen,
as a label attached after the fact --- characterizing what rest
produced, rather than steering how it got there. Whether this
descriptive label would be useful for anything downstream is not
examined here.

\paragraph{A check on construction, not a claim about meaning.} If a
later paper's pipeline is expected, empirically, to produce atoms
within some typical density/energy relationship, an atom that falls
far outside it --- for a free pair, there is no ``outside''; for a
bounded pair, violating the bound at all would be impossible by
construction and would indicate an error in how the atom was built,
not a new kind of atom. This is a debugging use, not an analytical
one, but worth keeping in mind precisely because a bounded pair
makes some configurations provably unreachable rather than merely
unlikely.

\section{Status}

The rule and its test are checked and stand on their own,
independent of any naming question. The four uses above are
genuinely speculative, offered as directions worth trying rather
than conclusions --- in particular, the Inference analogy needs
checking against that paper's own machinery before it is anything
more than a resemblance, and the debugging use depends on a bounded
pair actually being adopted into the formal inventory first.

\end{document}
